\documentclass[8pt,a4paper]{extarticle}
% ============================================================
% Gemeinsamer Preamble fuer alle 5 Open-Book-Spickzettel.
% Wird aus jedem klausur.tex via % ============================================================
% Gemeinsamer Preamble fuer alle 5 Open-Book-Spickzettel.
% Wird aus jedem klausur.tex via % ============================================================
% Gemeinsamer Preamble fuer alle 5 Open-Book-Spickzettel.
% Wird aus jedem klausur.tex via \input{../preamble.tex} geladen.
% Layout: A4 hochkant, 3 Spalten, 8pt, sehr kompakt.
% ============================================================

\usepackage[a4paper,margin=8mm,top=7mm,bottom=7mm,includefoot]{geometry}
\usepackage[utf8]{inputenc}
\usepackage[T1]{fontenc}
\usepackage[ngerman]{babel}
\usepackage{lmodern}
\usepackage{microtype}

\usepackage{amsmath,amssymb,amsfonts,mathtools}
\usepackage{multicol}
\usepackage{enumitem}
\usepackage{array,tabularx,booktabs,makecell}
\usepackage[dvipsnames,table]{xcolor}
\usepackage[explicit]{titlesec}
\usepackage{etoolbox}
\usepackage{fancyhdr}
\usepackage{lastpage}

% ---------- Farben ----------
\definecolor{secblue}{HTML}{1f4e79}
\definecolor{subblue}{HTML}{2e75b6}
\definecolor{boxbg}{HTML}{f2f2f2}
\definecolor{accent}{HTML}{c00000}
\definecolor{rulegray}{HTML}{888888}

% ---------- Spalten ----------
\setlength{\columnsep}{3mm}
\setlength{\columnseprule}{0.3pt}
\def\columnseprulecolor{\color{rulegray}}

% ---------- Listen kompakt ----------
\setlist{nosep,leftmargin=1.1em,labelsep=0.35em,topsep=0.2ex,partopsep=0pt}
\setlist[itemize]{label=\textbullet}

% ---------- Abstaende ----------
\setlength{\parindent}{0pt}
\setlength{\parskip}{0.15ex}
\setlength{\abovedisplayskip}{0.3ex}
\setlength{\belowdisplayskip}{0.3ex}
\setlength{\abovedisplayshortskip}{0.2ex}
\setlength{\belowdisplayshortskip}{0.2ex}
\setlength{\jot}{1pt}

% ---------- Tabellen kompakt ----------
\renewcommand{\arraystretch}{1.05}
\setlength{\tabcolsep}{3pt}
\setlength{\arrayrulewidth}{0.3pt}
\arrayrulecolor{rulegray}

% ---------- Section-Stil (kein Vollblock; dick + farbige Linie darunter) ----------
\titleformat{\section}
  {\large\bfseries\color{secblue}}
  {}{0pt}{#1}
  [\vspace{-0.4ex}{\color{secblue}\hrule height 0.8pt}\vspace{0.1ex}]
\titlespacing*{\section}{0pt}{1.4ex plus 0.4ex}{0.5ex}

\titleformat{\subsection}
  {\small\bfseries\color{subblue}}
  {}{0pt}{#1}
\titlespacing*{\subsection}{0pt}{0.9ex plus 0.2ex}{0.2ex}

\titleformat{\subsubsection}
  {\small\bfseries}
  {}{0pt}{#1}
\titlespacing*{\subsubsection}{0pt}{0.5ex}{0.1ex}

% ---------- Mini-Helfer ----------

% \fact{...} Hervorgehobene Regel (gerahmt, hellgrau)
\newcommand{\fact}[1]{%
  \begingroup\setlength{\fboxsep}{2pt}%
  \colorbox{boxbg}{\parbox{\dimexpr\linewidth-2\fboxsep}{\small #1\strut}}%
  \endgroup\par\vspace{0.2ex}}

% \warn{...} Stolperfalle (roter Strich links)
\newcommand{\warn}[1]{%
  \begingroup\setlength{\fboxsep}{2pt}%
  \noindent{\color{accent}\rule{2pt}{\baselineskip}}\hspace{2pt}%
  \begin{minipage}[t]{\dimexpr\linewidth-6pt}\small #1\end{minipage}%
  \endgroup\par\vspace{0.2ex}}

% Math-Shortcuts
\newcommand{\R}{\mathbb{R}}
\newcommand{\N}{\mathbb{N}}
\newcommand{\Z}{\mathbb{Z}}
\newcommand{\Q}{\mathbb{Q}}
\newcommand{\C}{\mathbb{C}}
\DeclareMathOperator{\rg}{rg}
\DeclareMathOperator{\spur}{tr}
\DeclareMathOperator{\im}{im}
\DeclareMathOperator{\diag}{diag}
\DeclareMathOperator{\sign}{sgn}
\DeclareMathOperator*{\argmin}{arg\,min}
\DeclareMathOperator*{\argmax}{arg\,max}
\newcommand{\trans}{^{\!\top}}
\newcommand{\norm}[1]{\left\lVert #1 \right\rVert}
\newcommand{\abs}[1]{\left\lvert #1 \right\rvert}
\newcommand{\inner}[2]{\langle #1,\,#2\rangle}
\newcommand{\dx}{\,\mathrm{d}x}
\newcommand{\dy}{\,\mathrm{d}y}
\newcommand{\dt}{\,\mathrm{d}t}

% Zeigt eine Display-Formel zentriert; skaliert nur runter wenn zu breit.
\newlength{\fitw}
\newcommand{\fitline}[1]{%
  \par\nointerlineskip
  \settowidth{\fitw}{$\displaystyle #1$}%
  \ifdim\fitw>\linewidth
    \noindent\resizebox{\linewidth}{!}{$\displaystyle #1$}\par
  \else
    \[\displaystyle #1\]%
  \fi}

% Tabular fuer fixe Spaltenbreiten (Spickzettel-typisch)
\newenvironment{tightt}[1]
  {\par\noindent\begin{tabular}{#1}}
  {\end{tabular}\par}
% Fuer X-Spalten direkt \begin{tabularx}{\linewidth}{...} verwenden
% (kann nicht in \newenvironment gewrappt werden -- TX@get@body Konflikt)

% Globale Schutzschalter gegen Overflow
\emergencystretch=2em
\tolerance=2000
\hbadness=10000
\hfuzz=2pt
\vfuzz=2pt

% ---------- Kopf-/Fusszeile ----------
\pagestyle{fancy}
\fancyhf{}
\renewcommand{\headrulewidth}{0pt}
\renewcommand{\footrulewidth}{0pt}
\fancyfoot[L]{\scriptsize\color{rulegray}\textit{\sheettitle}}
\fancyfoot[R]{\scriptsize\color{rulegray}Seite \thepage\,/\,\pageref{LastPage}}
\setlength{\footskip}{8pt}

% Titel-Header oben auf Seite 1
\newcommand{\sheetheader}[2]{%
  {\Large\bfseries\color{secblue} #1}\hfill
  {\small\color{rulegray} #2}\par
  \vspace{0.4ex}{\color{secblue}\hrule height 0.7pt}\vspace{1ex}}

% Default-Titel (wird ueberschrieben)
\newcommand{\sheettitle}{Open-Book Spickzettel}
 geladen.
% Layout: A4 hochkant, 3 Spalten, 8pt, sehr kompakt.
% ============================================================

\usepackage[a4paper,margin=8mm,top=7mm,bottom=7mm,includefoot]{geometry}
\usepackage[utf8]{inputenc}
\usepackage[T1]{fontenc}
\usepackage[ngerman]{babel}
\usepackage{lmodern}
\usepackage{microtype}

\usepackage{amsmath,amssymb,amsfonts,mathtools}
\usepackage{multicol}
\usepackage{enumitem}
\usepackage{array,tabularx,booktabs,makecell}
\usepackage[dvipsnames,table]{xcolor}
\usepackage[explicit]{titlesec}
\usepackage{etoolbox}
\usepackage{fancyhdr}
\usepackage{lastpage}

% ---------- Farben ----------
\definecolor{secblue}{HTML}{1f4e79}
\definecolor{subblue}{HTML}{2e75b6}
\definecolor{boxbg}{HTML}{f2f2f2}
\definecolor{accent}{HTML}{c00000}
\definecolor{rulegray}{HTML}{888888}

% ---------- Spalten ----------
\setlength{\columnsep}{3mm}
\setlength{\columnseprule}{0.3pt}
\def\columnseprulecolor{\color{rulegray}}

% ---------- Listen kompakt ----------
\setlist{nosep,leftmargin=1.1em,labelsep=0.35em,topsep=0.2ex,partopsep=0pt}
\setlist[itemize]{label=\textbullet}

% ---------- Abstaende ----------
\setlength{\parindent}{0pt}
\setlength{\parskip}{0.15ex}
\setlength{\abovedisplayskip}{0.3ex}
\setlength{\belowdisplayskip}{0.3ex}
\setlength{\abovedisplayshortskip}{0.2ex}
\setlength{\belowdisplayshortskip}{0.2ex}
\setlength{\jot}{1pt}

% ---------- Tabellen kompakt ----------
\renewcommand{\arraystretch}{1.05}
\setlength{\tabcolsep}{3pt}
\setlength{\arrayrulewidth}{0.3pt}
\arrayrulecolor{rulegray}

% ---------- Section-Stil (kein Vollblock; dick + farbige Linie darunter) ----------
\titleformat{\section}
  {\large\bfseries\color{secblue}}
  {}{0pt}{#1}
  [\vspace{-0.4ex}{\color{secblue}\hrule height 0.8pt}\vspace{0.1ex}]
\titlespacing*{\section}{0pt}{1.4ex plus 0.4ex}{0.5ex}

\titleformat{\subsection}
  {\small\bfseries\color{subblue}}
  {}{0pt}{#1}
\titlespacing*{\subsection}{0pt}{0.9ex plus 0.2ex}{0.2ex}

\titleformat{\subsubsection}
  {\small\bfseries}
  {}{0pt}{#1}
\titlespacing*{\subsubsection}{0pt}{0.5ex}{0.1ex}

% ---------- Mini-Helfer ----------

% \fact{...} Hervorgehobene Regel (gerahmt, hellgrau)
\newcommand{\fact}[1]{%
  \begingroup\setlength{\fboxsep}{2pt}%
  \colorbox{boxbg}{\parbox{\dimexpr\linewidth-2\fboxsep}{\small #1\strut}}%
  \endgroup\par\vspace{0.2ex}}

% \warn{...} Stolperfalle (roter Strich links)
\newcommand{\warn}[1]{%
  \begingroup\setlength{\fboxsep}{2pt}%
  \noindent{\color{accent}\rule{2pt}{\baselineskip}}\hspace{2pt}%
  \begin{minipage}[t]{\dimexpr\linewidth-6pt}\small #1\end{minipage}%
  \endgroup\par\vspace{0.2ex}}

% Math-Shortcuts
\newcommand{\R}{\mathbb{R}}
\newcommand{\N}{\mathbb{N}}
\newcommand{\Z}{\mathbb{Z}}
\newcommand{\Q}{\mathbb{Q}}
\newcommand{\C}{\mathbb{C}}
\DeclareMathOperator{\rg}{rg}
\DeclareMathOperator{\spur}{tr}
\DeclareMathOperator{\im}{im}
\DeclareMathOperator{\diag}{diag}
\DeclareMathOperator{\sign}{sgn}
\DeclareMathOperator*{\argmin}{arg\,min}
\DeclareMathOperator*{\argmax}{arg\,max}
\newcommand{\trans}{^{\!\top}}
\newcommand{\norm}[1]{\left\lVert #1 \right\rVert}
\newcommand{\abs}[1]{\left\lvert #1 \right\rvert}
\newcommand{\inner}[2]{\langle #1,\,#2\rangle}
\newcommand{\dx}{\,\mathrm{d}x}
\newcommand{\dy}{\,\mathrm{d}y}
\newcommand{\dt}{\,\mathrm{d}t}

% Zeigt eine Display-Formel zentriert; skaliert nur runter wenn zu breit.
\newlength{\fitw}
\newcommand{\fitline}[1]{%
  \par\nointerlineskip
  \settowidth{\fitw}{$\displaystyle #1$}%
  \ifdim\fitw>\linewidth
    \noindent\resizebox{\linewidth}{!}{$\displaystyle #1$}\par
  \else
    \[\displaystyle #1\]%
  \fi}

% Tabular fuer fixe Spaltenbreiten (Spickzettel-typisch)
\newenvironment{tightt}[1]
  {\par\noindent\begin{tabular}{#1}}
  {\end{tabular}\par}
% Fuer X-Spalten direkt \begin{tabularx}{\linewidth}{...} verwenden
% (kann nicht in \newenvironment gewrappt werden -- TX@get@body Konflikt)

% Globale Schutzschalter gegen Overflow
\emergencystretch=2em
\tolerance=2000
\hbadness=10000
\hfuzz=2pt
\vfuzz=2pt

% ---------- Kopf-/Fusszeile ----------
\pagestyle{fancy}
\fancyhf{}
\renewcommand{\headrulewidth}{0pt}
\renewcommand{\footrulewidth}{0pt}
\fancyfoot[L]{\scriptsize\color{rulegray}\textit{\sheettitle}}
\fancyfoot[R]{\scriptsize\color{rulegray}Seite \thepage\,/\,\pageref{LastPage}}
\setlength{\footskip}{8pt}

% Titel-Header oben auf Seite 1
\newcommand{\sheetheader}[2]{%
  {\Large\bfseries\color{secblue} #1}\hfill
  {\small\color{rulegray} #2}\par
  \vspace{0.4ex}{\color{secblue}\hrule height 0.7pt}\vspace{1ex}}

% Default-Titel (wird ueberschrieben)
\newcommand{\sheettitle}{Open-Book Spickzettel}
 geladen.
% Layout: A4 hochkant, 3 Spalten, 8pt, sehr kompakt.
% ============================================================

\usepackage[a4paper,margin=8mm,top=7mm,bottom=7mm,includefoot]{geometry}
\usepackage[utf8]{inputenc}
\usepackage[T1]{fontenc}
\usepackage[ngerman]{babel}
\usepackage{lmodern}
\usepackage{microtype}

\usepackage{amsmath,amssymb,amsfonts,mathtools}
\usepackage{multicol}
\usepackage{enumitem}
\usepackage{array,tabularx,booktabs,makecell}
\usepackage[dvipsnames,table]{xcolor}
\usepackage[explicit]{titlesec}
\usepackage{etoolbox}
\usepackage{fancyhdr}
\usepackage{lastpage}

% ---------- Farben ----------
\definecolor{secblue}{HTML}{1f4e79}
\definecolor{subblue}{HTML}{2e75b6}
\definecolor{boxbg}{HTML}{f2f2f2}
\definecolor{accent}{HTML}{c00000}
\definecolor{rulegray}{HTML}{888888}

% ---------- Spalten ----------
\setlength{\columnsep}{3mm}
\setlength{\columnseprule}{0.3pt}
\def\columnseprulecolor{\color{rulegray}}

% ---------- Listen kompakt ----------
\setlist{nosep,leftmargin=1.1em,labelsep=0.35em,topsep=0.2ex,partopsep=0pt}
\setlist[itemize]{label=\textbullet}

% ---------- Abstaende ----------
\setlength{\parindent}{0pt}
\setlength{\parskip}{0.15ex}
\setlength{\abovedisplayskip}{0.3ex}
\setlength{\belowdisplayskip}{0.3ex}
\setlength{\abovedisplayshortskip}{0.2ex}
\setlength{\belowdisplayshortskip}{0.2ex}
\setlength{\jot}{1pt}

% ---------- Tabellen kompakt ----------
\renewcommand{\arraystretch}{1.05}
\setlength{\tabcolsep}{3pt}
\setlength{\arrayrulewidth}{0.3pt}
\arrayrulecolor{rulegray}

% ---------- Section-Stil (kein Vollblock; dick + farbige Linie darunter) ----------
\titleformat{\section}
  {\large\bfseries\color{secblue}}
  {}{0pt}{#1}
  [\vspace{-0.4ex}{\color{secblue}\hrule height 0.8pt}\vspace{0.1ex}]
\titlespacing*{\section}{0pt}{1.4ex plus 0.4ex}{0.5ex}

\titleformat{\subsection}
  {\small\bfseries\color{subblue}}
  {}{0pt}{#1}
\titlespacing*{\subsection}{0pt}{0.9ex plus 0.2ex}{0.2ex}

\titleformat{\subsubsection}
  {\small\bfseries}
  {}{0pt}{#1}
\titlespacing*{\subsubsection}{0pt}{0.5ex}{0.1ex}

% ---------- Mini-Helfer ----------

% \fact{...} Hervorgehobene Regel (gerahmt, hellgrau)
\newcommand{\fact}[1]{%
  \begingroup\setlength{\fboxsep}{2pt}%
  \colorbox{boxbg}{\parbox{\dimexpr\linewidth-2\fboxsep}{\small #1\strut}}%
  \endgroup\par\vspace{0.2ex}}

% \warn{...} Stolperfalle (roter Strich links)
\newcommand{\warn}[1]{%
  \begingroup\setlength{\fboxsep}{2pt}%
  \noindent{\color{accent}\rule{2pt}{\baselineskip}}\hspace{2pt}%
  \begin{minipage}[t]{\dimexpr\linewidth-6pt}\small #1\end{minipage}%
  \endgroup\par\vspace{0.2ex}}

% Math-Shortcuts
\newcommand{\R}{\mathbb{R}}
\newcommand{\N}{\mathbb{N}}
\newcommand{\Z}{\mathbb{Z}}
\newcommand{\Q}{\mathbb{Q}}
\newcommand{\C}{\mathbb{C}}
\DeclareMathOperator{\rg}{rg}
\DeclareMathOperator{\spur}{tr}
\DeclareMathOperator{\im}{im}
\DeclareMathOperator{\diag}{diag}
\DeclareMathOperator{\sign}{sgn}
\DeclareMathOperator*{\argmin}{arg\,min}
\DeclareMathOperator*{\argmax}{arg\,max}
\newcommand{\trans}{^{\!\top}}
\newcommand{\norm}[1]{\left\lVert #1 \right\rVert}
\newcommand{\abs}[1]{\left\lvert #1 \right\rvert}
\newcommand{\inner}[2]{\langle #1,\,#2\rangle}
\newcommand{\dx}{\,\mathrm{d}x}
\newcommand{\dy}{\,\mathrm{d}y}
\newcommand{\dt}{\,\mathrm{d}t}

% Zeigt eine Display-Formel zentriert; skaliert nur runter wenn zu breit.
\newlength{\fitw}
\newcommand{\fitline}[1]{%
  \par\nointerlineskip
  \settowidth{\fitw}{$\displaystyle #1$}%
  \ifdim\fitw>\linewidth
    \noindent\resizebox{\linewidth}{!}{$\displaystyle #1$}\par
  \else
    \[\displaystyle #1\]%
  \fi}

% Tabular fuer fixe Spaltenbreiten (Spickzettel-typisch)
\newenvironment{tightt}[1]
  {\par\noindent\begin{tabular}{#1}}
  {\end{tabular}\par}
% Fuer X-Spalten direkt \begin{tabularx}{\linewidth}{...} verwenden
% (kann nicht in \newenvironment gewrappt werden -- TX@get@body Konflikt)

% Globale Schutzschalter gegen Overflow
\emergencystretch=2em
\tolerance=2000
\hbadness=10000
\hfuzz=2pt
\vfuzz=2pt

% ---------- Kopf-/Fusszeile ----------
\pagestyle{fancy}
\fancyhf{}
\renewcommand{\headrulewidth}{0pt}
\renewcommand{\footrulewidth}{0pt}
\fancyfoot[L]{\scriptsize\color{rulegray}\textit{\sheettitle}}
\fancyfoot[R]{\scriptsize\color{rulegray}Seite \thepage\,/\,\pageref{LastPage}}
\setlength{\footskip}{8pt}

% Titel-Header oben auf Seite 1
\newcommand{\sheetheader}[2]{%
  {\Large\bfseries\color{secblue} #1}\hfill
  {\small\color{rulegray} #2}\par
  \vspace{0.4ex}{\color{secblue}\hrule height 0.7pt}\vspace{1ex}}

% Default-Titel (wird ueberschrieben)
\newcommand{\sheettitle}{Open-Book Spickzettel}

\renewcommand{\sheettitle}{Optimierungsverfahren}

\begin{document}
\sheetheader{Optimierungsverfahren}{Open-Book Formelsammlung -- DHBW WDS125}

\begin{multicols*}{3}

% =================================================
\section{Grundlagen}
% =================================================
\subsection{Standardform}
\fitline{\min_{x\in X} f(x)\ \text{u.d.N.}\ g_i(x)\le 0,\ h_j(x)=0}
$\max f \equiv \min(-f)$ (Vorzeichen am Wert beachten).\\
Zul\"assige Menge $\mathcal{F}=\{x:g_i\le 0,\,h_j=0\}$.

\subsection{Klassifikation}
\begin{tabularx}{\linewidth}{@{}l@{\ }X@{}}
\textbf{LP} & $f,g,h$ linear -- Simplex, IPM \\
\textbf{QP} & $f$ quadr., $g/h$ linear -- Aktive-M. \\
\textbf{NLP} & mind.\ eines nichtlin.\ -- SQP, Newton \\
\textbf{Konvex} & lokal $=$ global \\
\textbf{Glatt} & $f\!\in\!C^1$ -- Gradientenverfahren \\
\textbf{Nichtglatt} & Subgrad./Proximal \\
\textbf{MIP/diskret} & $X\!\subseteq\!\Z^n$ -- B\&B \\
\textbf{Stochastisch} & SAA, Stoch.\ Gradient \\
\end{tabularx}

\subsection{Konvexit\"at}
\begin{itemize}
\item Menge: $\forall x,y\!\in\!C,\,\lambda\!\in\![0,1]\!:\,\lambda x{+}(1{-}\lambda)y\!\in\!C$
\item Funktion: $f(\lambda x{+}(1{-}\lambda)y)\le \lambda f(x){+}(1{-}\lambda)f(y)$
\item Strikt: $<$ f\"ur $x\ne y,\ \lambda\!\in\!(0,1)$
\item $f\in C^2$: konvex $\Leftrightarrow \nabla^2 f\succeq 0$ (PSD); strikt bei $\succ 0$
\item Summe konv., max endlich vieler konv., $f\!\circ\!\text{affin}$ konv.
\end{itemize}
\fact{\textbf{Hauptsatz}: konvex $\Rightarrow$ jedes lok.\ Min = glob.\ Min; strikt konvex $\Rightarrow$ eindeutig.}

\subsection{Optimalit\"at (unbeschr\"ankt)}
Notw.\ 1.\,Ord.: $\nabla f(x^\star)=0$ (station\"ar).\\
Hinr.\ 2.\,Ord.: $\nabla f(x^\star){=}0\,\wedge\,\nabla^2 f(x^\star)\succ 0\,\Rightarrow$ striktes lok.\ Min.\\
Konvex: $\nabla f(x^\star){=}0\,\Rightarrow$ glob.\ Min.

% =================================================
\section{Lineare Optimierung}
% =================================================
\subsection{Standardform LP}
\fitline{\min c\trans x\ \text{u.d.N.}\ Ax=b,\ x\ge 0}
\begin{itemize}
\item $\max c\trans x \to \min (-c\trans x)$
\item $a\trans x\le \beta \to a\trans x+s=\beta,\,s\ge 0$ (Schlupf)
\item $a\trans x\ge \beta \to a\trans x-s=\beta,\,s\ge 0$ (\"Uberschuss)
\item Freie Var.: $x=x^+-x^-$, $x^\pm\ge 0$
\end{itemize}
\fact{\textbf{Fundamentalsatz}: endliches Optimum $\Rightarrow$ optimale Ecke (Basisl\"osung) existiert.}

\subsection{Simplex}
\begin{enumerate}
\item \textbf{Opt.-Test}: alle red.\ Kosten $\bar c_j\!\ge\!0$? $\to$ fertig.
\item \textbf{Pivot-Spalte}: w\"ahle $j$ mit $\bar c_j\!<\!0$ (negativster = Dantzig).
\item \textbf{Pivot-Zeile} (min-ratio): minim.\ $b_i/a_{ij}$ unter $a_{ij}>0$; sonst LP unbeschr\"ankt.
\item \textbf{Pivot-Operation} um $a_{ij}$ (Zeile teilen, andere reduzieren).
\end{enumerate}
\warn{Bland-Regel (Tie $\to$ kleinster Index) verhindert Zyklen bei Degeneration. Worst case exponentiell (Klee-Minty), IPM polynomial.}

\subsection{Dualit\"at}
\textbf{Primal}: $\min c\trans x$, $Ax\ge b$, $x\ge 0$\\
\textbf{Dual}:\ $\max b\trans y$, $A\trans y\le c$, $y\ge 0$
\begin{tabularx}{\linewidth}{@{}l@{\ \,}X@{}}
$n$ Var (P) & $n$ NB (D) \\
$m$ NB (P) & $m$ Var (D) \\
$x_j\!\ge\!0$ & D-NB $j$ als $\le$ \\
$x_j$ frei & D-NB $j$ als $=$ \\
NB $i$ als $\ge$ & $y_i\!\ge\!0$ \\
NB $i$ als $=$ & $y_i$ frei \\
$c_j$ (Zielk.) & D-RHS \\
$b_i$ (RHS) & D-Zielk. \\
\end{tabularx}
\fact{\textbf{Schwach}: $b\trans y\le c\trans x$. \textbf{Stark}: endl.\ Opt.\ $\Rightarrow$ beide gleich.\\ \textbf{Kompl.\ Schlupf}: $x_j^\star(c_j{-}A_{\cdot j}\trans y^\star){=}0$,\ $y_i^\star(A_{i\cdot}x^\star{-}b_i){=}0$.}

% =================================================
\section{Nichtlineare Optimierung}
% =================================================
\subsection{1D-Verfahren}
\textbf{Bisektion auf $f'$}: $f'(a)f'(b)<0$, halbieren -- linear.\\
\textbf{Newton auf $f'$}: $x_{k+1}=x_k-f'(x_k)/f''(x_k)$ -- quadratisch.\\
\textbf{Goldener Schnitt}: $\rho{=}(\sqrt5{-}1)/2\!\approx\!0{,}618$;\ $x_1{=}a{+}(1{-}\rho)(b{-}a)$,\ $x_2{=}a{+}\rho(b{-}a)$;\ kleineres $f$ behalten.

\subsection{Gradientenabstieg}
\fitline{x_{k+1}=x_k-\alpha_k\,\nabla f(x_k)}
\textbf{Armijo}: $f(x_k{-}\alpha\nabla f)\le f(x_k){-}c_1\alpha\|\nabla f\|^2$, $c_1\!\in\!(0,\tfrac12)$.\\
\textbf{Wolfe}: Armijo + Kr\"ummungsbed.\ \\
Konvergenzrate (Kondition $\kappa$): $((\kappa{-}1)/(\kappa{+}1))^2$ -- langsam bei $\kappa\!\gg\!1$.

\subsection{Newton \& Quasi-Newton}
\fitline{x_{k+1}=x_k-\bigl[\nabla^2 f(x_k)\bigr]^{-1}\nabla f(x_k)}
\begin{itemize}
\item Quadratisch lokal; Hesse muss PD f\"ur Abstieg.
\item Aufwand $O(n^3)$/Schritt; ged\"ampft mit Linesearch.
\item \textbf{Levenberg-Marquardt}: $(\nabla^2 f+\lambda I)^{-1}$.
\item \textbf{BFGS}: approx.\ $H^{-1}$ aus Gradientendiff., superlinear.
\end{itemize}

\subsection{Lagrange \& KKT}
\fitline{L(x,\lambda,\mu)=f(x)+\sum_i \lambda_i g_i(x)+\sum_j \mu_j h_j(x)}
$\lambda_i\ge 0$ (Ungl.), $\mu_j$ frei (Gl.).\\[0.3ex]
\textbf{KKT} (notwendig bei reg.\ Punkt; bei konvex auch hinreichend):
\begin{enumerate}
\item Station.: $\nabla f+\sum_i\lambda_i\nabla g_i+\sum_j\mu_j\nabla h_j=0$
\item Primal-Zul.: $g_i\le 0,\ h_j=0$
\item Dual-Zul.: $\lambda_i\ge 0$
\item Kompl.\ Schlupf: $\lambda_i\cdot g_i=0$
\end{enumerate}

\subsection{Strafverfahren}
\textbf{Quadr.\ Penalty}: $f+\rho\sum h_j^2+\rho\sum\max(0,g_i)^2$, $\rho\!\to\!\infty$.\\
\textbf{Barriere/IPM}: $f-\mu\sum\log(-g_i)$, $\mu\!\to\!0$.

% =================================================
\section{Gradientenfreie / Metaheuristiken}
% =================================================
\subsection{Direktsuche}
\textbf{Nelder-Mead}: Simplex aus $n{+}1$ Punkten; Schritte: Spiegelung $\to$ Expansion (wenn besser als alle) / Kontraktion / Schrumpfung.\\
\textbf{Koordinatenabstieg}: zyklisch eine Variable opt.\\
\textbf{Hooke-Jeeves}: Achsen-Schritte + Pattern Move.

\subsection{Simulated Annealing}
\begin{enumerate}
\item Initial $x_0,\ T_0$; best$=x_0$.
\item Wiederhole bis $T<T_{\min}$:
  \begin{itemize}
  \item Nachbar $x'$; $\Delta=f(x'){-}f(x)$.
  \item Akzeptiere wenn $\Delta<0$ oder $\text{rand}()<\exp(-\Delta/T)$.
  \item Best updaten.
  \end{itemize}
\item Cooling $T_{k+1}=\alpha T_k$, $\alpha\!\in\![0{,}9;\,0{,}99]$.
\end{enumerate}
\fact{Metropolis: $P_{\text{accept}}=\exp(-\Delta/T)$. Log.\ Cooling $T_k=c/\log(k{+}1)$ konv.\ zum Glob.\ (theor.).}

\subsection{Genetische Algorithmen}
Selektion (Roulette/Turnier/Rang) -- Crossover (One-Point/Uniform/BLX-$\alpha$) -- Mutation (Gauss/Bit-Flip) -- Ersetzung $(\mu,\lambda)$ vs.\ $(\mu{+}\lambda)$; \textbf{Elitismus} bewahrt Beste.\\
\textbf{CMA-ES}: adaptive Kovarianzmatrix.

\subsection{PSO}
\fitline{v_i \leftarrow w\,v_i+c_1 r_1(p_i-x_i)+c_2 r_2(g-x_i)}
\fitline{x_i \leftarrow x_i+v_i}
$w\!\in\![0{,}4;\,0{,}9]$ (linear fallend), $c_1{=}c_2\!\approx\!2$, $r_{1,2}\!\sim\!U[0,1]$.\ $p_i$ pers., $g$ glob.\ Best.

\subsection{ACO}
$p_{ij}\propto \tau_{ij}^\alpha\cdot\eta_{ij}^\beta$, $\eta_{ij}{=}1/d_{ij}$, $\tau$ verstaerkt/verdunstet.

% =================================================
\section{Mehrzieloptimierung}
% =================================================
\subsection{Pareto}
\textbf{Dominanz} $x\prec y$: $\forall i\!:f_i(x)\!\le\!f_i(y)\,\wedge\,\exists j\!:f_j(x)\!<\!f_j(y)$.\\
Pareto-optimal $=$ nicht dominiert; Pareto-Front $=$ Bild der Pareto-Menge.

\subsection{Verfahren}
\textbf{Gewichtete Summe}: $\min\sum w_i f_i$, $w_i\!\ge\!0$, $\sum w_i{=}1$ -- bei nicht-konvexer Front nicht alle Punkte erreichbar.\\
\textbf{$\epsilon$-Constraint}: $\min f_1$ u.d.N.\ $f_j\!\le\!\epsilon_j$ ($j\ne 1$) -- auch nicht-konvex.\\
\textbf{NSGA-II / SPEA2 / MOEA/D}.

\subsection{Qualit\"atsindikatoren}
\textbf{Hypervolumen} (h\"oher besser, Pareto-konform); \textbf{IGD} (niedriger besser); \textbf{Spread}.

% =================================================
\section{Verfahrenswahl-Matrix}
% =================================================
\begin{tabularx}{\linewidth}{@{}l@{\ \,}X@{}}
LP & Simplex / IPM \\
QP & Aktive-M., IPM \\
NLP glatt & SQP, Newton, BFGS \\
Konvex & Newton/BFGS \\
Nicht-glatt & Subgrad., Proximal \\
MIP/diskret & B\&B, Heuristik \\
Blackbox/verrauscht & Nelder-Mead, CMA-ES, Bayes \\
Nicht-konvex & Multi-Start, SA, GA, PSO \\
Mehrzielig & Gew.\ Sum., $\epsilon$-C, NSGA-II \\
$n\!<\!100$ & Dense Newton, SLSQP/IPOPT \\
$n\!<\!10^4$ & BFGS, sparse Newton, CG \\
$n\!>\!10^5$ (DL) & L-BFGS, SGD/Adam \\
\end{tabularx}

% =================================================
\section{Klausur-Stolperfallen}
% =================================================
\begin{itemize}
\item $\lambda_i\ge 0$ nur bei Ungl.-NB; $\mu_j$ frei bei Gleichungs-NB.
\item Komplement.\ Schlupf bei KKT \emph{immer} mitschreiben.
\item Min-Ratio nur \"uber $a_{ij}>0$ (sonst LP unbeschr\"ankt).
\item Bland-Regel bei Degeneration aktivieren.
\item Bei $\max$ Vorzeichenwechsel: Optimum $-(-c)\trans x^\star$.
\item Gewichtete Summe verfehlt nicht-konvexe Pareto-Fronten.
\item Vor SA/GA: Problem konvex? dann Newton/BFGS.
\end{itemize}

\end{multicols*}
\end{document}
